% ============================================================
%  GUÍA DE ESTUDIO: INTEGRACIÓN POR CAMBIO DE VARIABLE Y PARTES
%  Autor: Ing. Gabriel Astudillo
%  Estructura según estructura.md
% ============================================================

\documentclass[12pt, a4paper]{article}

% ── Paquetes ─────────────────────────────────────────────────

\usepackage{mathwizards-guia}
\mwguia{Cambio de Variable y Por Partes}{Ing. Gabriel Astudillo $\cdot$ Guía de Estudio}

% ── Comandos propios de esta guía ───────────────────
\newcommand{\dx}{\,dx}
\newcommand{\du}{\,du}
\newcommand{\dt}{\,dt}
\newcommand{\dv}{\,dv}
\newcommand{\dw}{\,dw}
\newcommand{\dtheta}{\,d\theta}
\newcommand{\R}{\mathbb{R}}
\newcommand{\CC}{\,+\,C}

\begin{document}
% ============================================================

% ── Portada / Título ─────────────────────────────────────────
\mwportada{Guía de Estudio y Ejercicios}{Integración por Cambio de Variable y Por Partes}{Teoría, Métodos y 70 Problemas Progresivos con Respuestas}

\vspace{4pt}
\begin{cuadroinfo}
\small
\textbf{Presentación:} Esta guía está diseñada para llevarte desde los conceptos fundamentales de la integración no directa hasta la resolución de problemas desafiantes de cálculo integral.
Consta de explicaciones conceptuales detalladas basadas en las reglas de derivación, estrategias paso a paso para la selección de variables y 70 ejercicios organizados progresivamente por nivel de complejidad, finalizando con su respectiva clave de respuestas.
\\[4pt]
\textbf{Nivel de Dificultad:}\quad
\facil{} Básico / Directo \quad
\medio{} Intermedio \quad
\dificil{} Desafiante / Avanzado
\end{cuadroinfo}

\vspace{8pt}

% ============================================================
\section{El Método de Cambio de Variable (Sustitución)}
% ============================================================

\subsection{\textquestiondown Qué es un cambio de variable?}
El método de integración por \textbf{cambio de variable} (o sustitución) es la herramienta fundamental que nos permite deshacer la \textbf{Regla de la Cadena} de la derivación.

Recordemos que si tenemos una función compuesta $y = F(g(x))$, su derivada respecto a $x$ se expresa como:
$$\frac{d}{dx} \big[ F(g(x)) \big] = F'(g(x)) \cdot g'(x)$$

Si integramos en ambos lados con respecto a $x$, obtenemos:
$$\int F'(g(x)) \cdot g'(x) \dx = F(g(x)) \CC$$

Si definimos una nueva variable independiente $u = g(x)$, su diferencial asociado será $du = g'(x)\dx$. Al sustituir directamente estos términos en la integral original, la transformamos en una expresión mucho más sencilla de resolver:
\begin{cuadroteoria}
$$\int f(g(x)) \cdot g'(x) \dx = \int f(u) \du = F(u) \CC$$
Donde $f(u) = F'(u)$.
\end{cuadroteoria}
Este proceso reduce una integral compleja en una integral de tipo elemental o directa sobre la variable auxiliar $u$.

\subsection{\textquestiondown Cuándo usar el cambio de variable?}
Debes considerar el cambio de variable cuando en la integral reconozcas:
\begin{enumerate}[label=\alph*)]
  \item Una \textbf{función compuesta} $f(g(x))$ donde la derivada del argumento interno, $g'(x)$ (o un factor proporcional a ella), se encuentra multiplicando al resto del integrando.
  \item Una estructura algebraica fraccionaria o irracional complicada donde al realizar una asignación del tipo $u = g(x)$ se logre linealizar o simplificar el integrando, permitiendo cancelar términos en $x$ mediante despejes directos.
\end{enumerate}

\subsection{\textquestiondown Cómo elegir el cambio de variable correcto?}
La elección del cambio de variable óptimo suele seguir estas recomendaciones prácticas:
\begin{cuadrosolucion}
\begin{itemize}[leftmargin=*, label=\checkmark]
  \item \textbf{El argumento interno:} Si tienes expresiones como $\sin(g(x))$, $e^{g(x)}$, o $\sqrt{g(x)}$, intenta con $u = g(x)$.
  \item \textbf{El denominador:} En funciones racionales $\int \frac{h(x)}{g(x)} \dx$, prueba haciendo $u = g(x)$ (el denominador completo), ya que su derivada podría estar en el numerador.
  \item \textbf{La base de una potencia:} Si observas $[g(x)]^n$, asigna $u = g(x)$.
\end{itemize}
\end{cuadrosolucion}

\noindent\textbf{Procedimiento general en 5 pasos:}
\begin{enumerate}
  \item \textbf{Identificar} y proponer la sustitución $u = g(x)$.
  \item \textbf{Calcular el diferencial} $du = g'(x) \dx$. Despejar $\dx$ si te resulta más cómodo para sustituir.
  \item \textbf{Reescribir la integral} en términos de $u$. Asegúrate de que no quede ninguna variable $x$. Si quedan términos en $x$, intenta despejar $x$ a partir de la ecuación original $u = g(x)$.
  \item \textbf{Resolver} la integral resultante con respecto a $u$.
  \item \textbf{Volver a la variable original} reemplazarizando $u$ por su definición original $g(x)$.
\end{enumerate}

\noindent\textbf{Ejemplo Guiado 1 (Sustitución Directa):} Evaluar $\int x^2 e^{x^3} \dx$.
\begin{itemize}
  \item Identificamos la función interna $g(x) = x^3$, cuya derivada es $3x^2$ (proporcional al factor $x^2$ del integrando).
  \item Definimos $u = x^3 \implies du = 3x^2 \dx \implies x^2 \dx = \frac{1}{3} \du$.
  \item Reemplazamos en la integral:
    $$\int x^2 e^{x^3} \dx = \int e^{u} \left(\frac{1}{3}\du\right) = \frac{1}{3}\int e^u \du = \frac{1}{3}e^u \CC$$
  \item Volviendo a la variable $x$: $\frac{1}{3}e^{x^3} \CC$.
\end{itemize}

\noindent\textbf{Ejemplo Guiado 2 (Sustitución con Despeje):} Evaluar $\int x\sqrt{x-1}\dx$.
\begin{itemize}
  \item Definimos $u = x-1 \implies du = \dx$.
  \item Notamos que queda una $x$ libre. La despejamos a partir del cambio: $x = u+1$.
  \item Sustituimos todo en la integral:
    $$\int x\sqrt{x-1}\dx = \int (u+1)\sqrt{u}\du = \int (u^{3/2} + u^{1/2})\du$$
  \item Integramos: $\frac{2}{5}u^{5/2} + \frac{2}{3}u^{3/2} \CC$.
  \item Volvemos a $x$: $\frac{2}{5}(x-1)^{5/2} + \frac{2}{3}(x-1)^{3/2} \CC$.
\end{itemize}

\newpage

% ============================================================
\subsection{Ejercicios Progresivos: Cambio de Variable}
% ============================================================
Resuelve las siguientes integrales indefinidas utilizando el método de cambio de variable.

\begin{multicols}{2}
\begin{enumerate}[leftmargin=*]
  \item \facil{} $\int (3x - 5)^7 \dx$
  \item \facil{} $\int \cos(4x + 1) \dx$
  \item \facil{} $\int e^{-2x} \dx$
  \item \facil{} $\int \frac{1}{2x + 7} \dx$
  \item \facil{} $\int \sec^2(5x) \dx$
  \item \medio{} $\int 2x(x^2 + 3)^5 \dx$
  \item \medio{} $\int x^2 e^{x^3} \dx$
  \item \medio{} $\int \frac{\ln x}{x} \dx$
  \item \medio{} $\int \sin^4 x \cos x \dx$
  \item \medio{} $\int \frac{x}{\sqrt{x^2 + 9}} \dx$
  \item \medio{} $\int \frac{e^x}{1 + e^{2x}} \dx$
  \item \medio{} $\int \frac{\sec^2 x}{\tan x} \dx$
  \item \medio{} $\int \cos x \sqrt{\sin x} \dx$
  \item \dificil{} $\int x \sqrt{x - 1} \dx$
  \item \dificil{} $\int \frac{x^2}{x - 2} \dx$
  \item \dificil{} $\int \frac{\sin(1/x)}{x^2} \dx$
  \item \dificil{} $\int e^{\cos x} \sin x \dx$
  \item \dificil{} $\int \frac{x^3}{\sqrt{x^2 + 1}} \dx$
  \item \dificil{} $\int \frac{1}{x \ln x \ln(\ln x)} \dx$
  \item \dificil{} $\int \frac{\cos(\ln x)}{x} \dx$
\end{enumerate}
\end{multicols}

\vspace{10pt}

% ============================================================
\section{El Método de Integración por Partes}
% ============================================================

\subsection{\textquestiondown Qué es la integración por partes?}
El método de \textbf{integración por partes} permite integrar funciones que se expresan como el producto de dos términos algebraicos o trascendentes que no están relacionados por sus derivadas directas. Nace de revertir la \textbf{Regla del Producto} de la diferenciación.

Si tenemos dos funciones derivables $u(x)$ y $v(x)$, la regla para la derivada de su producto es:
$$\frac{d}{dx} \big[ u(x)v(x) \big] = u'(x)v(x) + u(x)v'(x)$$

Al integrar ambos miembros de esta ecuación con respecto a $x$:
$$u(x)v(x) = \int u'(x)v(x) \dx + \int u(x)v'(x) \dx$$

Reorganizando e introduciendo los diferenciales $du = u'(x)\dx$ y $dv = v'(x)\dx$, obtenemos la célebre fórmula de integración por partes:
\begin{cuadroadvertencia}
$$\int u \dv = u \cdot v - \int v \du$$
\end{cuadroadvertencia}

\subsection{\textquestiondown Cuándo usar la integración por partes?}
Este método es ideal cuando el integrando presenta alguna de las siguientes características:
\begin{itemize}
  \item Es el producto de dos funciones de distinta clase (por ejemplo, el producto de un polinomio por una función exponencial o trigonométrica).
  \item Consiste en una única función cuya derivada conocemos perfectamente, pero cuya antiderivada directa no está tabulada (como $\ln x$, $\arctan x$, $\arcsin x$, etc.).
\end{itemize}

\subsection{\textquestiondown Cómo elegir la integración por partes correcta?}
Para asignar con éxito las expresiones correspondientes a $u$ y a $dv$, se suele emplear la regla mnemotécnica \textbf{ILATE}, la cual indica el orden de prioridad para elegir la función $u$ de izquierda a derecha:
\begin{cuadroinfo}
\begin{description}
  \item[\textbf{I} - Inversas trigonométricas] ($\arcsin x$, $\arccos x$, $\arctan x$, \dots)
  \item[\textbf{L} - Logarítmicas] ($\ln x$, $\log_a x$, \dots)
  \item[\textbf{A} - Algebráicas / Polinomiales] ($x^n$, $\sqrt{x}$, $x^2 + 3$, \dots)
  \item[\textbf{T} - Trigonométricas] ($\sin x$, $\cos x$, $\sec x$, \dots)
  \item[\textbf{E} - Exponenciales] ($e^x$, $2^x$, \dots)
\end{description}
\end{cuadroinfo}
La función elegida como $u$ será diferenciada para obtener $du$, mientras que el resto del integrando (que obligatoriamente incluye al diferencial $dx$) se define como $dv$ y debe ser integrada para hallar $v$.

\noindent\textbf{Ejemplo Guiado:} Evaluar $\int x \cos x \dx$.
\begin{itemize}
  \item Analizamos el integrando: $x$ es una función algebraica (A), y $\cos x$ es trigonométrica (T). Según \textbf{ILATE}, la algebraica tiene prioridad para $u$.
  \item Definimos:
    \begin{align*}
      u &= x \implies du = \dx \\
      dv &= \cos x \dx \implies v = \sin x
    \end{align*}
  \item Aplicamos la fórmula:
    $$\int x\cos x\dx = x\sin x - \int \sin x \dx$$
  \item Resolvemos la integral restante:
    $$\int x\cos x\dx = x\sin x - (-\cos x) \CC = x\sin x + \cos x \CC$$
\end{itemize}

\newpage

% ============================================================
\subsection{Ejercicios Progresivos: Integración por Partes}
% ============================================================
Resuelve las siguientes integrales indefinidas utilizando el método de integración por partes. En este bloque no se incluyen integrales que requieran partes sucesivas o cíclicas, aunque algunas pueden requerir un cambio de variable auxiliar.

\begin{multicols}{2}
\begin{enumerate}[leftmargin=*, resume]
  \item \facil{} $\int x e^x \dx$
  \item \facil{} $\int x \cos x \dx$
  \item \facil{} $\int x \sin(2x) \dx$
  \item \facil{} $\int (2x - 3) e^{-x} \dx$
  \item \medio{} $\int x \sec^2 x \dx$
  \item \medio{} $\int (x + 2) \sin(3x) \dx$
  \item \medio{} $\int x 3^x \dx$
  \item \medio{} $\int (5x + 1) \cos(4x) \dx$
  \item \medio{} $\int \ln x \dx$
  \item \medio{} $\int x \ln x \dx$
  \item \medio{} $\int \arcsin x \dx$
  \item \medio{} $\int \arctan x \dx$
  \item \medio{} $\int x^2 \ln x \dx$
  \item \dificil{} $\int \frac{\ln x}{x^2} \dx$
  \item \dificil{} $\int x \arctan x \dx$
  \item \dificil{} $\int \ln(2x + 1) \dx$
  \item \dificil{} $\int e^{\sqrt{x}} \dx$
  \item \dificil{} $\int x^3 e^{x^2} \dx$
  \item \dificil{} $\int \frac{x \arcsin x}{\sqrt{1-x^2}} \dx$
  \item \dificil{} $\int \cos(\sqrt{x}) \dx$
\end{enumerate}
\end{multicols}

\vspace{10pt}

% ============================================================
\section{Integración por Partes Sucesiva y Cíclicas}
% ============================================================

\subsection{Definiciones y Fundamentos}
Al aplicar el método por partes, es frecuente toparse con escenarios que requieren consideraciones especiales:

\begin{enumerate}[leftmargin=*]
  \item \textbf{Integración por Partes Sucesiva:}
    Ocurre cuando el polinomio $u = x^n$ tiene un grado $n \ge 2$. Para poder anular el polinomio a través de sus sucesivas derivadas, debemos aplicar la fórmula por partes de manera reiterada $n$ veces consecutivas.
    
  \item \textbf{Integración por Partes Cíclica:}
    Se presenta cuando integramos productos de funciones trascendentes cuyas derivadas son periódicas (como exponenciales combinadas con senos o cosenos, o la integral de $\sec^3 x$). Al aplicar el método dos veces, volvemos a obtener un término equivalente a la integral original con signo opuesto. En lugar de continuar infinitamente, resolvemos el problema de forma \textbf{algebraica}, despejando la integral de la ecuación resultante.
\end{enumerate}

\noindent\textbf{Ejemplo Guiado (Partes Sucesivas):} Evaluar $\int x^2 e^x \dx$.
\begin{itemize}
  \item Tomamos $u_1 = x^2 \implies du_1 = 2x\dx$, y $dv_1 = e^x\dx \implies v_1 = e^x$.
  \item Primera aplicación:
    $$\int x^2 e^x \dx = x^2 e^x - 2\int x e^x \dx$$
  \item Para la nueva integral $\int x e^x \dx$, aplicamos partes nuevamente:
    $$u_2 = x \implies du_2 = \dx, \quad dv_2 = e^x\dx \implies v_2 = e^x$$
  \item Segunda aplicación:
    $$\int x^2 e^x \dx = x^2 e^x - 2\left(x e^x - \int e^x \dx\right) = x^2 e^x - 2x e^x + 2e^x \CC$$
  \item Factorizando: $(x^2 - 2x + 2)e^x \CC$.
\end{itemize}

\noindent\textbf{Ejemplo Guiado (Partes Cíclicas):} Evaluar $I = \int e^x \sin x \dx$.
\begin{itemize}
  \item Elegimos $u_1 = \sin x \implies du_1 = \cos x \dx$, y $dv_1 = e^x \dx \implies v_1 = e^x$.
  \item Primera aplicación:
    $$I = e^x \sin x - \int e^x \cos x \dx$$
  \item Aplicamos partes a $\int e^x \cos x \dx$ con $u_2 = \cos x \implies du_2 = -\sin x \dx$, y $dv_2 = e^x \dx \implies v_2 = e^x$:
    $$\int e^x \cos x \dx = e^x \cos x - \int e^x (-\sin x)\dx = e^x \cos x + \int e^x \sin x \dx$$
  \item Sustituimos de vuelta en $I$:
    $$I = e^x \sin x - \left(e^x \cos x + I\right) = e^x \sin x - e^x \cos x - I$$
  \item Sumamos $I$ a ambos lados de la ecuación:
    $$2I = e^x(\sin x - \cos x) \implies I = \frac{1}{2}e^x(\sin x - \cos x) \CC$$
\end{itemize}

\newpage

% ============================================================
\subsection{Ejercicios Progresivos: Sucesivas y Cíclicas}
% ============================================================
Resuelve las siguientes integrales indefinidas utilizando integración por partes sucesiva o cíclica.

\begin{multicols}{2}
\begin{enumerate}[leftmargin=*, resume]
  \item \facil{} $\int x^2 e^x \dx$
  \item \facil{} $\int x^2 \sin x \dx$
  \item \facil{} $\int x^2 \cos(2x) \dx$
  \item \medio{} $\int x^3 e^{-x} \dx$
  \item \medio{} $\int (\ln x)^2 \dx$
  \item \medio{} $\int x(\ln x)^2 \dx$
  \item \medio{} $\int x^2 e^{3x} \dx$
  \item \medio{} $\int x^3 \cos x \dx$
  \item \medio{} $\int e^x \sin x \dx$
  \item \medio{} $\int e^x \cos x \dx$
  \item \dificil{} $\int e^{2x} \sin(3x) \dx$
  \item \dificil{} $\int e^{-x} \cos(2x) \dx$
  \item \dificil{} $\int \sin(\ln x) \dx$
  \item \dificil{} $\int \cos(\ln x) \dx$
  \item \dificil{} $\int \sec^3 x \dx$
\end{enumerate}
\end{multicols}

\vspace{10pt}

% ============================================================
\section{Combinación de Técnicas de Integración}
% ============================================================

\subsection{Estrategia General y Ejemplos}
En la práctica, las integrales reales rara vez vienen clasificadas con la técnica que deben resolverse. A menudo requieren combinar varios métodos para su resolución definitiva.
Para enfrentarte a estas situaciones, te sugerimos seguir esta hoja de ruta mental:

\begin{cuadroextra}
\begin{enumerate}[leftmargin=*, label=\arabic*.]
  \item \textbf{Simplificación Algebraica:} \textquestiondown Puedo reescribir la expresión utilizando propiedades algebraicas o identidades trigonométricas? (por ejemplo, expandir binomios, dividir polinomios, identidades de ángulos dobles).
  \item \textbf{Sustitución Inicial:} \textquestiondown Existe alguna función interna molesta o un argumento no lineal (como $\sqrt{x}$, $e^x$, $\ln x$, $\arctan x$) cuya eliminación mediante cambio de variable reduzca la integral a una forma conocida?
  \item \textbf{Por Partes:} Si después de sustituir nos queda un producto de funciones de diferente índole, aplicamos la regla de integración por partes.
\end{enumerate}
\end{cuadroextra}

\noindent\textbf{Ejemplo Guiado:} Evaluar $\int \frac{\arctan e^x}{e^x} \dx$.
\begin{itemize}
  \item Vemos que el argumento de la arcotangente es $e^x$ y hay términos exponenciales en el denominador.
  \item Aplicamos un cambio de variable inicial: $u = e^x \implies du = e^x \dx \implies \dx = \frac{\du}{u}$.
  \item Sustituimos en la integral:
    $$\int \frac{\arctan e^x}{e^x} \dx = \int \frac{\arctan u}{u} \cdot \frac{\du}{u} = \int \frac{\arctan u}{u^2} \du = \int u^{-2} \arctan u \du$$
  \item Esta integral requiere ahora el método de integración por partes:
    \begin{align*}
      y &= \arctan u \implies dy = \frac{1}{1+u^2} \du \\
      dz &= u^{-2} \du \implies z = -\frac{1}{u}
    \end{align*}
  \item Aplicamos partes:
    $$\int u^{-2} \arctan u \du = -\frac{\arctan u}{u} + \int \frac{1}{u(1+u^2)} \du$$
  \item Resolvemos la integral racional restante por fracciones parciales:
    $$\frac{1}{u(1+u^2)} = \frac{1}{u} - \frac{u}{1+u^2}$$
    $$\int \left(\frac{1}{u} - \frac{u}{1+u^2}\right) \du = \ln|u| - \frac{1}{2}\ln(1+u^2) \CC$$
  \item Reagrupamos y volvemos a la variable $x$ original sabiendo que $u = e^x$ e $1+u^2 = 1+e^{2x}$:
    $$\int \frac{\arctan e^x}{e^x} \dx = -\frac{\arctan e^x}{e^x} + \ln(e^x) - \frac{1}{2}\ln(1+e^{2x}) \CC$$
    $$= x - e^{-x}\arctan e^x - \frac{1}{2}\ln(1+e^{2x}) \CC$$
\end{itemize}

\newpage

% ============================================================
\subsection{Ejercicios Progresivos: Combinación de Técnicas}
% ============================================================
Resuelve las siguientes integrales utilizando combinaciones de las técnicas estudiadas (simplificación trigonométrica, cambio de variable, integración por partes y álgebra elemental).

\begin{multicols}{2}
\begin{enumerate}[leftmargin=*, resume]
  \item \medio{} $\int \frac{e^{2x}}{\sqrt{e^x + 1}} \dx$
  \item \medio{} $\int \ln(x^2 + 1) \dx$
  \item \medio{} $\int \frac{\ln(\ln x)}{x} \dx$
  \item \medio{} $\int \frac{x^5}{x^3 + 1} \dx$
  \item \medio{} $\int \frac{\arcsin(\sqrt{x})}{\sqrt{x}} \dx$
  \item \medio{} $\int x^5 e^{x^3} \dx$
  \item \medio{} $\int \cos^3 x \sqrt{\sin x} \dx$
  \item \medio{} $\int x^2 \arctan x \dx$
  \item \dificil{} $\int \frac{x e^{\arctan x}}{(1+x^2)^{3/2}} \dx$
  \item \dificil{} $\int \frac{\ln(x + \sqrt{x^2+1})}{\sqrt{x^2+1}} \dx$
  \item \dificil{} $\int x^3 \sin(x^2) \dx$
  \item \dificil{} $\int \frac{\arctan e^x}{e^x} \dx$
  \item \dificil{} $\int \sec^3 x \tan^3 x \dx$
  \item \dificil{} $\int \frac{x \ln(x^2+1)}{x^2+1} \dx$
  \item \dificil{} $\int \frac{e^x(x-1)}{x^2} \dx$
\end{enumerate}
\end{multicols}

\vspace{10pt}

% ============================================================
\section{Resumen de Criterios y Consejos Útiles}
% ============================================================
\begin{table}[h!]
\centering
\small
\begin{tabular}{p{3.5cm} p{5.5cm} p{6.5cm}}
\toprule
\textbf{Método} & \textbf{\textquestiondown Cuándo aplicar?} & \textbf{Estrategia / Tip clave} \\
\midrule
\textbf{Cambio de Variable (Sustitución)} & Integrales que contienen una función interna $g(x)$ y su derivada $g'(x)$ multiplicando. & Elige $u = g(x)$. El diferencial $du$ debe absorber el resto del integrando. \\
\midrule
\textbf{Integración por Partes (ILATE)} & Producto de funciones de distintas familias (algebraicas, trigonométricas, exponenciales, logarítmicas). & Usa la regla mnemotécnica \textbf{ILATE} para elegir $u$. La parte más difícil que se pueda integrar pasa a ser $dv$. \\
\midrule
\textbf{Partes Sucesivas} & Polinomios multiplicados por funciones que no se extinguen en la integral ($\sin x$, $e^x$). & Deriva reiteradamente el término algebraico hasta que su grado se reduzca a cero. \\
\midrule
\textbf{Partes Cíclicas} & Combinaciones de tipo $e^{ax}\cos(bx)$ o potencias impares de funciones trigonométricas. & Aplica el método 2 veces hasta reaparecer la integral original, y despéjala algebraicamente. \\
\midrule
\textbf{Combinación} & Expresiones con argumentos complejos y coeficientes de diferente índole. & Aplica primero una sustitución para simplificar argumentos y luego resuelve el producto resultante por partes. \\
\bottomrule
\end{tabular}
\caption{Criterio de selección de técnicas para la resolución de integrales.}
\label{tab:resumen}
\end{table}

\newpage

% ============================================================
\ifmwclaves
\section{Clave de Respuestas}
% ============================================================
A continuación se presenta la solución analítica para cada uno de los 70 ejercicios planteados en esta guía. Todas las respuestas incluyen la constante de integración $C \in \R$.

\begin{enumerate}[leftmargin=*, label=\textbf{\arabic*.}]
  \item $\frac{1}{24}(3x-5)^8 \CC$
  \item $\frac{1}{4}\sin(4x+1) \CC$
  \item $-\frac{1}{2}e^{-2x} \CC$
  \item $\frac{1}{2}\ln|2x+7| \CC$
  \item $\frac{1}{5}\tan(5x) \CC$
  \item $\frac{1}{6}(x^2+3)^6 \CC$
  \item $\frac{1}{3}e^{x^3} \CC$
  \item $\frac{1}{2}(\ln x)^2 \CC$
  \item $\frac{1}{5}\sin^5 x \CC$
  \item $\sqrt{x^2+9} \CC$
  \item $\arctan(e^x) \CC$
  \item $\ln|\tan x| \CC$
  \item $\frac{2}{3}(\sin x)^{3/2} \CC$
  \item $\frac{2}{15}(x-1)^{3/2}(3x+2) \CC$
  \item $\frac{1}{2}x^2 + 2x + 4\ln|x-2| \CC$
  \item $\cos\left(\frac{1}{x}\right) \CC$
  \item $-e^{\cos x} \CC$
  \item $\frac{1}{3}\sqrt{x^2+1}(x^2-2) \CC$
  \item $\ln|\ln(\ln x)| \CC$
  \item $\sin(\ln x) \CC$
  \item $(x-1)e^x \CC$
  \item $x\sin x + \cos x \CC$
  \item $-\frac{1}{2}x\cos(2x) + \frac{1}{4}\sin(2x) \CC$
  \item $(1-2x)e^{-x} \CC$
  \item $x\tan x + \ln|\cos x| \CC$
  \item $-\frac{1}{3}(x+2)\cos(3x) + \frac{1}{9}\sin(3x) \CC$
  \item $\frac{3^x}{\ln 3}\left(x - \frac{1}{\ln 3}\right) \CC$
  \item $\frac{1}{4}(5x+1)\sin(4x) + \frac{5}{16}\cos(4x) \CC$
  \item $x\ln x - x \CC$
  \item $\frac{1}{2}x^2\ln x - \frac{1}{4}x^2 \CC$
  \item $x\arcsin x + \sqrt{1-x^2} \CC$
  \item $x\arctan x - \frac{1}{2}\ln(1+x^2) \CC$
  \item $\frac{1}{3}x^3\ln x - \frac{1}{9}x^3 \CC$
  \item $-\frac{\ln x + 1}{x} \CC$
  \item $\frac{1}{2}(x^2+1)\arctan x - \frac{1}{2}x \CC$
  \item $\left(x+\frac{1}{2}\right)\ln(2x+1) - x \CC$
  \item $2(\sqrt{x}-1)e^{\sqrt{x}} \CC$
  \item $\frac{1}{2}(x^2-1)e^{x^2} \CC$
  \item $x - \sqrt{1-x^2}\arcsin x \CC$
  \item $2\sqrt{x}\sin(\sqrt{x}) + 2\cos(\sqrt{x}) \CC$
  \item $(x^2 - 2x + 2)e^x \CC$
  \item $-x^2\cos x + 2x\sin x + 2\cos x \CC$
  \item $\frac{1}{2}x\cos(2x) + \frac{1}{4}(2x^2-1)\sin(2x) \CC$
  \item $-(x^3 + 3x^2 + 6x + 6)e^{-x} \CC$
  \item $x(\ln x)^2 - 2x\ln x + 2x \CC$
  \item $\frac{1}{4}x^2\left[2(\ln x)^2 - 2\ln x + 1\right] \CC$
  \item $\frac{1}{27}(9x^2 - 6x + 2)e^{3x} \CC$
  \item $(x^3 - 6x)\sin x + (3x^2 - 6)\cos x \CC$
  \item $\frac{1}{2}e^x(\sin x - \cos x) \CC$
  \item $\frac{1}{2}e^x(\sin x + \cos x) \CC$
  \item $\frac{1}{13}e^{2x}(2\sin(3x) - 3\cos(3x)) \CC$
  \item $\frac{1}{5}e^{-x}(2\sin(2x) - \cos(2x)) \CC$
  \item $\frac{x}{2}(\sin(\ln x) - \cos(\ln x)) \CC$
  \item $\frac{x}{2}(\sin(\ln x) + \cos(\ln x)) \CC$
  \item $\frac{1}{2}\sec x\tan x + \frac{1}{2}\ln|\sec x + \tan x| \CC$
  \item $\frac{2}{3}(e^x-2)\sqrt{e^x+1} \CC$
  \item $x\ln(x^2+1) - 2x + 2\arctan x \CC$
  \item $\ln x[\ln(\ln x) - 1] \CC$
  \item $\frac{1}{3}x^3 - \frac{1}{3}\ln(x^3+1) \CC$
  \item $2\sqrt{x}\arcsin(\sqrt{x}) + 2\sqrt{1-x} \CC$
  \item $\frac{1}{3}(x^3-1)e^{x^3} \CC$
  \item $\frac{2}{3}(\sin x)^{3/2} - \frac{2}{7}(\sin x)^{7/2} \CC$
  \item $\frac{1}{3}x^3\arctan x - \frac{1}{6}x^2 + \frac{1}{6}\ln(1+x^2) \CC$
  \item $\frac{(x-1)e^{\arctan x}}{2\sqrt{1+x^2}} \CC$
  \item $\frac{1}{2}[\ln(x + \sqrt{x^2+1})]^2 \CC$
  \item $\frac{1}{2}\sin(x^2) - \frac{1}{2}x^2\cos(x^2) \CC$
  \item $x - e^{-x}\arctan e^x - \frac{1}{2}\ln(1+e^{2x}) \CC$
  \item $\frac{1}{5}\sec^5 x - \frac{1}{3}\sec^3 x \CC$
  \item $\frac{1}{4}[\ln(x^2+1)]^2 \CC$
  \item $\frac{e^x}{x} \CC$
\end{enumerate}

\fi
\end{document}
